\documentclass[11pt,a4paper]{article}

\usepackage[a4paper,margin=1in]{geometry}
\usepackage{fontspec}
\IfFontExistsTF{Times New Roman}{\setmainfont{Times New Roman}}{\setmainfont{TeX Gyre Termes}}
\IfFontExistsTF{Menlo}{\setmonofont{Menlo}}{\setmonofont{Latin Modern Mono}}
\usepackage{amsmath,amssymb}
\usepackage{booktabs}
\usepackage{longtable}
\usepackage{array}
\usepackage{enumitem}
\usepackage{graphicx}
\usepackage{fancyhdr}
\usepackage[round,authoryear]{natbib}
\usepackage[hidelinks]{hyperref}
\usepackage{xurl}
\usepackage{microtype}

\setlength{\parindent}{0pt}
\setlength{\parskip}{0.65em}
\setlist[itemize]{leftmargin=1.4em}
\newcommand{\file}[1]{\path{#1}}
\newcommand{\code}[1]{\path{#1}}
\newcolumntype{L}[1]{>{\raggedright\arraybackslash}p{#1}}

\pagestyle{fancy}
\fancyhf{}
\lhead{\small v1 Small-Grid Dataset}
\rhead{\small Dataset QA Review}
\cfoot{\thepage}
\renewcommand{\headrulewidth}{0.4pt}

\title{\textbf{v1 Small-Grid Dataset Report}\\
\large 288-Regime Surrogate Dataset Generation and QA}
\author{Codex-assisted surrogate data review}
\date{June 18, 2026}

\begin{document}
\maketitle

\begin{abstract}
This report reviews the Stage 2 v1 small-grid dataset generation for the American option
risk-surface project. The generation used exactly the 288 approved regimes from the construction
plan, the validated baseline American Crank--Nicolson / projected SOR solver, and the established
LCP, boundary, and Greek diagnostic layers. It produced one compressed \code{.npz} package,
human-audit manifests, QA summaries, QA figures, and a higher-grid confirmation table. It did not
train a neural model, create \code{.pt} files, switch to Rannacher smoothing, or generate regimes
outside the approved v1 plan. The v1 review decision is
\code{READY_FOR_PRICE_SURROGATE_PLANNING}.
\end{abstract}

\tableofcontents
\newpage

\section{Purpose of v1 Small-Grid Dataset}

The v1 small-grid dataset is the first approved full small-grid dataset package after the v0
dry-run generation test. Its purpose is to create a compact, auditable dataset for future
price-surrogate planning. The dataset stores normalized features, price and continuation-premium
labels, diagnostic boundary and Greek fields, masks, solver metadata, split assignments, and QA
manifests.

This is not a production dataset and not a model-training stage. It is a small-grid dataset and
QA gate that comes before any neural surrogate experiment.

\section{Roadmap Link and v0 Result}

The compressed roadmap places this work in Stage 2: v1 small-grid dataset generation plus dataset
QA. The preceding v0 dry run accepted all eight dry-run regimes and produced 41,624 accepted
sampled rows. That result supported moving from pipeline testing to the approved 288-regime
small-grid generation.

The solver source remains:
\begin{center}
\code{solver_name = american_crank_nicolson_psor_price}\\
\code{solver_variant = baseline_cn_psor}
\end{center}
The numerical foundation follows the validated American finite-difference and LCP framework used
throughout the solver ladder \citep{black_scholes_1973,merton_1973,brennan_schwartz_1977,wilmott_howison_dewynne_1995}.

\section{Regime Count and Split Counts}

The v1 generation used the exact Cartesian product approved in the construction plan:
\[
\{\mathrm{put},\mathrm{call}\}\times
\{0.25,0.5,1.0,2.0\}\times
\{0.20,0.40,0.60\}\times
\{0.01,0.05,0.10\}\times
\{0.00,0.03,0.06,0.10\}.
\]
This gives 288 regimes: 144 puts and 144 calls. No-dividend call regimes remain explicit control
regimes rather than dividend-call early-exercise evidence.

\begin{longtable}{L{0.22\textwidth} L{0.20\textwidth} L{0.20\textwidth} L{0.28\textwidth}}
\toprule
Split & Regimes & Accepted rows & Purpose \\
\midrule
\endhead
Train & 202 & 1,051,006 & Main fitting split for future price-surrogate planning. \\
Validation & 19 & 98,857 & Regime-level validation split. \\
Test & 43 & 223,729 & Regime-level test and control regimes. \\
Stress holdout & 24 & 124,872 & Held-out stress combinations. \\
\midrule
Total & 288 & 1,498,464 & All accepted v1 sampled rows. \\
\bottomrule
\end{longtable}

\section{Solver and Grid Settings}

Every v1 regime used \(K=1\), \(S_{\max}=4\), \(M=N=120\), and the baseline CN/PSOR solver.
Sampling was restricted to the approved reporting and learning region \(S/K\in[0.4,1.8]\) and
\(\tau/T\in[0,1]\). The full solver domain was retained internally for boundary conditions and
diagnostics.

Each accepted regime contributed 5,203 sampled rows: 43 spot nodes by 121 time rows. The
compressed package size is 50,262,900 bytes.

\section{Feature, Label, and Mask Schema}

The feature matrix uses:
\[
x=\left(\log(S/K),\ \tau/T,\ r,\ q,\ \sigma,\ T,\ \mathrm{is\_call}\right).
\]
The stored labels are:
\begin{itemize}
  \item core labels: \code{value_over_K}, \code{payoff_over_K}, and \code{premium_over_K};
  \item diagnostic labels: \code{exercise_indicator} and \code{boundary_spot_over_K};
  \item cautioned Greek diagnostics: \code{delta} and \code{scaled_gamma}.
\end{itemize}

The mask matrix contains payoff-kink-near, boundary-near, maturity-row, strict-interior,
Gamma-allowed, Delta-allowed, exercise-region, and continuation-region masks. These masks are
not optional: they are the main protection against treating boundary and Gamma diagnostics as
production-quality labels.

\section{Diagnostics and Acceptance Checks}

A regime was accepted only if PSOR converged, obstacle violation was at most \(10^{-8}\), equation
violation was at most \(10^{-6}\), complementarity product was at most \(10^{-6}\), and required
metadata and masks were present.

\begin{longtable}{L{0.38\textwidth} L{0.24\textwidth} L{0.28\textwidth}}
\toprule
Metric & v1 result & Interpretation \\
\midrule
\endhead
Regimes solved & 288 & Exactly the approved v1 plan. \\
Accepted regimes & 288 & All regimes passed acceptance checks. \\
Accepted sampled rows & 1,498,464 & 5,203 rows per accepted regime. \\
Maximum obstacle violation & \(0.0\) & Payoff obstacle respected in accepted regimes. \\
Maximum equation violation & \(1.02\times 10^{-7}\) & Below the \(10^{-6}\) threshold. \\
Maximum complementarity product & \(1.82\times 10^{-8}\) & Below the \(10^{-6}\) threshold. \\
Boundary-found time rows & 29,795 & Boundary metadata was recorded across regimes. \\
Strict Greek mask nodes & 3,825,842 & Greek diagnostics remain mask-dependent. \\
\bottomrule
\end{longtable}

\section{Output Files and Manifests}

\begin{longtable}{L{0.42\textwidth} L{0.13\textwidth} L{0.33\textwidth}}
\toprule
Output & Type & Purpose \\
\midrule
\endhead
\file{results/04_surrogate_dataset/v1_small_grid/dataset_v1_small_grid.npz} & \code{.npz} & Accepted-row v1 arrays only. \\
\file{results/04_surrogate_dataset/v1_small_grid/regime_manifest.csv} & CSV & Per-regime parameters, diagnostics, and acceptance status. \\
\file{results/04_surrogate_dataset/v1_small_grid/diagnostic_summary.csv} & CSV & PSOR, LCP, boundary, Greek, and runtime diagnostics. \\
\file{results/04_surrogate_dataset/v1_small_grid/split_assignment.csv} & CSV & Regime-level split assignment and row counts. \\
\file{results/04_surrogate_dataset/v1_small_grid/schema_snapshot.csv} & CSV & Feature, label, mask, and audit-field schema. \\
\file{results/04_surrogate_dataset/v1_small_grid/output_manifest.csv} & CSV & Output audit manifest for v1 artifacts. \\
\file{results/04_surrogate_dataset/v1_small_grid/generation_log.txt} & TXT & Plain-text generation progress log. \\
\file{results/04_surrogate_dataset/v1_small_grid/v1_dataset_quality_summary.csv} & CSV & Aggregate QA status and decision metrics. \\
\file{results/04_surrogate_dataset/v1_small_grid/v1_higher_grid_confirmation.csv} & CSV & QA-only \(M=N=180\) confirmation checks. \\
\bottomrule
\end{longtable}

The compressed package contains the keys \code{X}, \code{y_value}, \code{y_payoff},
\code{y_premium}, \code{y_exercise_indicator}, \code{y_boundary}, \code{y_delta},
\code{y_scaled_gamma}, \code{masks}, \code{regime_index}, \code{feature_names},
\code{label_names}, \code{mask_names}, \code{audit_numeric}, \code{audit_numeric_names},
\code{regime_ids}, and \code{split_names}.

\section{QA Figures}

The QA figures are diagnostic only. They check row balance, diagnostic tolerances, premium
distribution, boundary metadata availability, and Greek-mask availability.

\begin{figure}[htbp]
\centering
\includegraphics[width=0.82\textwidth]{results/04_surrogate_dataset/v1_small_grid/figures/v1_split_row_counts.png}
\caption{Accepted sampled rows by regime-level split. The counts match the deterministic construction plan.}
\end{figure}

\begin{figure}[htbp]
\centering
\includegraphics[width=0.82\textwidth]{results/04_surrogate_dataset/v1_small_grid/figures/v1_diagnostic_thresholds.png}
\caption{Obstacle, equation, and complementarity diagnostics by regime. The maxima remain below the v1 acceptance thresholds.}
\end{figure}

\begin{figure}[htbp]
\centering
\includegraphics[width=0.78\textwidth]{results/04_surrogate_dataset/v1_small_grid/figures/v1_premium_distribution.png}
\caption{Continuation-premium distribution over accepted rows. This supports price and premium label QA, not production risk claims.}
\end{figure}

\begin{figure}[htbp]
\centering
\includegraphics[width=0.70\textwidth]{results/04_surrogate_dataset/v1_small_grid/figures/v1_boundary_availability.png}
\caption{Mean boundary-found time rows by option family. Boundary values remain threshold-based diagnostics.}
\end{figure}

\begin{figure}[htbp]
\centering
\includegraphics[width=0.70\textwidth]{results/04_surrogate_dataset/v1_small_grid/figures/v1_gamma_mask_summary.png}
\caption{Greek mask node counts by option family. Gamma remains a cautioned diagnostic field and should be used only with masks.}
\end{figure}

\section{Higher-Grid Confirmation}

Four representative regimes were rerun at \(M=N=180\) for QA only. These confirmation rows were
not mixed into the v1 dataset arrays. All four confirmation runs passed the same diagnostic
acceptance rule. The maximum selected-node value difference versus the v1 grid was about
\(1.11\times 10^{-2}\), with mean absolute value differences near \(1.98\times 10^{-3}\).
This confirms basic scale consistency while preserving the v1 dataset as an \(M=N=120\) package.

\section{v1 Dataset Review Decision}

The v1 review decision is:
\begin{center}
\fbox{\code{READY_FOR_PRICE_SURROGATE_PLANNING}}
\end{center}

This means the small-grid package is ready to support a future price-surrogate planning step. It
does not automatically approve neural training, production deployment, or unmasked Greek labels.

\section{Limitations}

\begin{itemize}
  \item The dataset is a small-grid package, not a production-scale dataset.
  \item Boundary labels remain threshold-based diagnostics, not exact free-boundary truth.
  \item Delta and scaled Gamma remain diagnostic fields. Gamma is high caution near payoff kinks,
  maturity rows, and extracted boundary bands.
  \item The \(M=N=180\) confirmation is small and QA-oriented; it is not a full convergence proof.
  \item No neural model has been trained or validated from this dataset.
\end{itemize}

\section{What This Dataset Supports}

The v1 package supports price-surrogate planning with normalized features, reliable price and
premium labels, split metadata, diagnostic boundary fields, Greek masks, solver diagnostics, and
regime-level manifests. It also supports human review of whether the first surrogate experiment
should target price and continuation premium before using any Greek diagnostic targets.

\section{What This Dataset Does Not Support}

The v1 package does not support production risk claims, final paper conclusions by itself,
production Greek labels, broad extrapolation to regimes outside the approved grid, or automatic
approval to train neural surrogates without a separate surrogate-planning gate.

\section{Recommended Next Step}

The recommended next step is a price-surrogate planning gate. That gate should decide the first
model family, training objective, loss functions, regime-level evaluation protocol, and whether
Greek fields remain metadata-only for the first surrogate experiment.

\bibliographystyle{plainnat}
\bibliography{reports/03_solver/references}

\end{document}
